\chapter{Challenges of Social Transformation: Women, Caste and Ethnicity}

\section{Atrocities Against Women: Religion and Structure}

\begin{KeyConcept}
\begin{itemize}
\item \textbf{Religion is one basis of atrocities against women} --- in every religion there are structures enabling men (and society) to dominate women:
\item \textbf{Hinduism} --- women are given \textbf{Shudra status} (Manusmriti), cannot be high priests or light the cremation fire; \textbf{menstrual blood is seen as polluting} (the Sabarimala exclusion, despite the Supreme Court); the \textbf{Kul Badhu} (a woman is responsible for the whole family); the idealised images of \textbf{Sati and Savitri}; and the \textbf{husband as demi-god} (bhakti). Veena Das: woman is at once \textbf{Devi and Dasi}.
\item \textbf{Islam} --- \textbf{Patricia Uberoi}: Muslim marriage is a \textbf{contract} (with \textbf{maher}), but the maher is usually small and in over 80\% of cases \textbf{not paid}; the \textbf{burqa} is a patriarchal construct (socialised as 'natural', tied to family dignity).
\item \textbf{Catholicism} --- divorce and abortion are seen as \textbf{sin}, so a married woman becomes a 'prisoner of the family'.
\item The atrocities are of \textbf{three types}: \textbf{social} (female feticide/infanticide, dowry, child marriage, child labour), \textbf{domestic} (rape, sexual atrocities, domestic violence, mental/physical torture), and \textbf{violence} (manifest and latent).
\item \textbf{Manifest violence} = realised (rape, physical torture, eve-teasing, stalking). \textbf{Latent violence} = unrealised/structural (the sexual division of labour; self-inflicted violence --- women not eating, not taking care of themselves, the feminisation of poverty, nutrition going last to women).
\item \textbf{Rape is an exhibition of power}, not biological need: \textbf{social power} (an upper-caste man raping a Dalit woman = caste + patriarchy), \textbf{economic power} (pink-collar jobs, the feminisation of agriculture), and \textbf{political power} (no representation; the women's reservation bill, Mulayam Singh's opposition, Nitish Kumar's women's reservation). \textbf{Renuka Singh}: men control women's bodies in \textbf{legitimate (marriage)} and \textbf{illegitimate (rape)} ways.
\item \textbf{Nivedita Menon}: India legalised abortion early --- but for \textbf{population control} (a neo-liberal idea), not feminism; the \textbf{PCPNDT Act} (sex-determination) came the same way.
\end{itemize}
\end{KeyConcept}

\begin{KeyConcept}
\begin{itemize}
\item \textbf{Social structure} (Levi-Strauss's binary opposites) is part of our \textbf{unconscious} --- so women are also patriarchal (socialised into it). \textbf{Gender is performative} (Judith Butler) and \textbf{socially constructed} (Simone de Beauvoir: 'one is not born, but rather becomes, a woman'). \textbf{Constructivism}: caste, gender, even 'child' are socially constructed.
\item \textbf{Sandhya Srinivasan}: women cooking with traditional fuels for $\sim$3 hours/day inhale smoke equal to \textbf{20 cigarettes} --- an unintended (latent) violence.
\item \textbf{Ajantha Subramanian} ('Cast of Merit'): IITs are socialised in a casteist/Brahminical way --- ISRO succeeds but nobody built a technology for \textbf{manual scavengers} (the Bandicoot robot, made by IIT students, never got funding); Abdul Kalam (missiles) is celebrated, small innovations are not.
\item \textbf{Utsa Patnaik}: women work \textbf{3,980 hours/year} vs men's 1,230 (a pair of bullocks: 1,084) --- worse than cattle.
\item \textbf{Bina Agarwal}: machines removed women from agriculture (men sit on machines), and men create \textbf{myths} (that machines kill, that women are weak, that engineering is not for women) to keep women away from technology --- leading to \textbf{domestication} and \textbf{pink-collar jobs} (nurse, receptionist, teacher --- the 'feminine' jobs in the least productive sector).
\end{itemize}
\end{KeyConcept}

\begin{QuickRevision}
\begin{itemize}
\item Religion as basis: Hindu (Shudra status, menstrual pollution/Sabarimala, Kul Badhu, Sati/Savitri, husband as demi-god), Islam (maher not paid, burqa), Catholic (divorce/abortion = sin).
\item Three types: social, domestic, violence; latent (structural, self-inflicted) vs manifest (rape, torture).
\item Rape = exhibition of power (social/economic/political); Renuka Singh (legitimate/illegitimate control); Nivedita Menon (abortion = population control; PCPNDT).
\item Structure (binary opposites, unconscious); gender = performative/constructed; Srinivasan (20 cigarettes), Subramanian (Cast of Merit), Patnaik (3980 hrs), Agarwal (machines/myths/pink-collar).
\end{itemize}
\end{QuickRevision}

\section{Domestic Violence}

\begin{KeyConcept}
\begin{itemize}
\item Domestic violence was only \textbf{criminalised in 1983} (Section 498A IPC --- cruelty by husband/in-laws, cognizable and non-bailable); the \textbf{Domestic Violence Act} followed (2002/2005). Before that, the state never saw it as violence --- it was treated as a \textbf{private matter} (others should not intervene).
\item \textbf{Neera Desai} (TISS report): as long as there is \textbf{tolerance for domestic violence in cultural, legal and political institutions}, laws alone cannot work --- and laws made \emph{within} the system only reinforce the (patriarchal) hierarchies.
\item \textbf{Uma Chakravarti} ('The Frog in the Well'): after marriage the father's and mother's roles are opposite --- the father says 'if there is a problem, phone', the mother says 'adjust'.
\item \textbf{Reasons}: the grip of family and patriarchy; \textbf{dowry} (mental/physical torture); not bearing a (male) child; \textbf{earning more than the husband} (urban India). Women are seen not as autonomous but \textbf{in relation to male family members}.
\item \textbf{Class difference}: upper/middle-class women (no childhood experience of violence) suffer \textbf{psychological trauma} and cannot fight back; lower-class women (who have lived among violence) \textbf{fight back and talk about it} --- and have more freedom (they are not domesticated).
\end{itemize}
\end{KeyConcept}

\begin{QuickRevision}
\begin{itemize}
\item 1983 (S.498A) $\rightarrow$ domestic violence criminalised; DV Act later; treated as private matter.
\item Neera Desai (TISS): tolerance in institutions; laws within the system reinforce hierarchy.
\item Uma Chakravarti ('Frog in the Well'): father 'phone', mother 'adjust'; reasons: dowry, no child, earning more; women not autonomous; class difference in response.
\end{itemize}
\end{QuickRevision}

\section{Caste Conflict}

\begin{KeyConcept}
\begin{itemize}
\item Caste conflict resurged after the \textbf{1970s}. It began with the British (equality before law $\rightarrow$ empowerment), initially as \textbf{upper-caste vs Dalit} conflict; after independence the \textbf{dominant (backward) castes} (land-owners --- Jats, Yadavs, Kurmis in the North) took the space vacated by the upper castes and \textbf{enforced the traditional social order} (Tanjore, Tamil Nadu, 1990s --- a Dalit basti burned, $\sim$100 dead; Ranveer Sena vs MCC in Bihar over land).
\item The \textbf{Dalit Panthers} (1970s, Marx + Ambedkar) produced \textbf{Dalit literature} (Dalits explaining their own consciousness) and the \textbf{SC/ST Act}.
\item \textbf{Kanshi Ram}: '\textbf{Bahujan}' (multiple heads); \textbf{Mayawati}: '\textbf{Sarvajan}' (a Brahmin--Dalit coalition --- the exploiter and exploited against another exploiter); \textbf{Bhim Army} (education $\rightarrow$ social mobility); \textbf{Bhima Koregaon} (asserting that Dalits \emph{have} histories --- exclusion from history is itself violence); \textbf{Jignesh Mevani}; \textbf{Gopal Guru} ('subaltern Hinduism' --- Hindutva bringing Dalits into its fold). \textbf{Pauline Kolenda}: the sweepers' own myth (they were Brahmins, outcasted) --- caste groups create mythologies to resist the karma logic.
\item \textbf{Muzzaffar Assadi and S. Rajendra} --- causes of caste conflict: \textbf{land issues} (Dalits who bought land couldn't farm it --- no irrigation, tubewells; joined Naxal groups), \textbf{temple entry}, \textbf{forcing human excreta} (dry latrines; banned by a 1994 act), \textbf{nude worship} (Karnataka), \textbf{grazing-ground conflict}, \textbf{bonded labour and social boycott}, \textbf{untouchability in hotels}, \textbf{discrimination in bureaucracy/workplace}, and \textbf{false implication in cases} (the film \emph{Aakrosh}).
\end{itemize}
\end{KeyConcept}

\begin{QuickRevision}
\begin{itemize}
\item Caste conflict resurgent after 1970s; British equality $\rightarrow$ upper-caste vs Dalit; dominant castes enforced order (Tanjore, Ranveer Sena/MCC).
\item Dalit Panthers (Dalit literature + SC/ST Act); Kanshi Ram (Bahujan), Mayawati (Sarvajan), Bhim Army, Bhima Koregaon, Gopal Guru (subaltern Hinduism).
\item Assadi \& Rajendra: land, temple entry, human excreta/dry latrines (1994 act), nude worship, grazing, bonded labour, discrimination, false implication (Aakrosh).
\end{itemize}
\end{QuickRevision}

\section{Ethnic Conflict and Ethnicity}

\begin{KeyConcept}
\begin{itemize}
\item \textbf{Ethnicity} can be defined \textbf{narrowly} (racial/linguistic groups --- a colonial legacy of dividing coloniser from colonised) or \textbf{broadly} (\textbf{Horowitz}): 'ethnicity designates a sense of collective belonging based on common descent, language, history, culture, race or religion, or some combination' --- so Protestant--Catholic, Hindu--Muslim, black--white, Tamil--Sinhala and Shia--Sunni are \textbf{all ethnic conflicts} in the broad sense (religion-based conflict is called 'communal' in India).
\item \textbf{Ethnic group vs nation}: an ethnic group \textbf{may do without a state}, but a \textbf{nation implies bringing ethnicity and statehood together}; \textbf{Ernst Gellner}: nationalism is 'the principle that the political and the national unit should be congruent'.
\item \textbf{Causes of ethnic conflict} (four theories/perceptions): \textbf{primordial identity} (ethnic bonds are stronger than civic/national bonds --- the idea of 'ancient hatred'), the role of \textbf{myths} (myths create charismatic personalities and legitimise actions --- e.g. the mythology around Gandhi as 'Mahatma', or Modi's 'Bal Narendra' crocodile story; and Sanskritisation through fake genealogies --- e.g. the Yadavs of Yadavpur appointing a Brahmin to write a genealogy).
\end{itemize}
\end{KeyConcept}

\begin{QuickRevision}
\begin{itemize}
\item Ethnicity: narrow (race/language) vs broad (Horowitz --- collective belonging by descent/language/history/culture/race/religion).
\item Ethnic group (no state) vs nation (ethnicity + statehood); Gellner (political = national unit).
\item Causes: primordial identity/'ancient hatred'; myths (charisma + legitimation; Yadavpur fake genealogy).
\end{itemize}
\end{QuickRevision}
